\chapter{第三讲补编}

此结果的一个重要推论如下:
\begin{theorem}\label{thm:6.1}
对任意渐近平直初数据集 $(\Sigma,\bar g,\bar k)$, 其极大未来发展 $(\mathcal{M},g)$ 必可找到一个具有紧闭包的区域 $\Omega_0\subset\Sigma$, 使得其在 $\mathcal{M}$ 中影响域 $\mathcal{C}^+(\Omega_0)$ 的边界 $\mathcal{D}_l^+$ 拥有完备的零测地线生成元.\end{theorem}

上述推论可用来引入"完备未来零无穷远"的概念:
\begin{definition}\label{def:6.2}
若对任一正实数 $A$, 我们都能找到一个包含定理 1.2 中区域 $\Omega_0$ 的区域 $\Omega$, 使得 $\Omega$ 在 $\mathcal{M}$ 中依赖域的边界 $\mathcal{D}^-$具有如下性质: 它的每条零测地线生成元, 从 $\mathcal{D}^-\cap\mathcal{D}_0^+$ 起算的总仿射长度至少为 $A$, 则称渐近平直初数据集 $(\Sigma,\bar g,\bar k)$ 的极大未来发展 $(\mathcal{M},g)$ 拥有\textbf{完备未来零无穷远}.\end{definition}

对于充分大的初数据集, 奇点的不可避免出现, 促使 Penrose 提出两个猜想, 即弱宇宙监督猜想与强宇宙监督猜想.
\noindent\textbf{猜想 1(弱宇宙监督猜想 WCC)}(\textit{Conjecture 1, Weak Cosmic Censorship})．\label{conj:wcc}一般的渐近平直初数据, 其极大未来发展拥有完备的未来零无穷远.
\par\medskip

WCC 猜想并不排除奇点可能被局部观测者看到的可能性. 这可能导致如下悖论式的情形: 此类观测者所作观测的结果缺乏唯一可预测性. 强宇宙监督猜想正是为了禁止奇点这一不良特性而提出的.
\noindent\textbf{猜想 2(强宇宙监督猜想 SCC)}(\textit{Conjecture 2, Strong Cosmic Censorship})．\label{conj:scc}一般的初数据集, 其极大未来发展作为 Lorentz 流形, 在连续意义下局部不可延拓.
\par\medskip

更技术地说, 这意味着除去某些可能的例外初值条件外, 一个初数据集的极大未来发展具有如下性质: 沿任一未来、不可延拓、有限长度的类时测地线, 时空曲率分量(相对沿该测地线平行移动的标准正交标架表示)当弧长趋近于其极限值时必趋于无穷.
\section*{6.3 \thmref{thm:3.10} 到\thmref{thm:3.12} 的化简}

下面我们描述\thmref{thm:3.10}(以下记为定理 (A))到\thmref{thm:3.12} 之衰减估计(以下记为定理 (B))证明的化简. 前五个步骤如今已是标准的, 参见 \cite{klai2}, 尤其是 \cite{klai72}; 最后一步的化简则是 \cite{klai29} 几何方法的典型特征.
\section*{6.4 步骤 1: 能量估计}

如前所述, 我们对拟线性波动方程的能量估计表明: 只要 $\|\partial\phi\|_{L_{[0,T]}^1L_x^\infty}\leq 1$, 解 $\phi$ 的 $L_{[0,T]}^\infty H^s$ 范数即被初数据 $\phi[0]$ 的 $H^s$ 范数控制. 下面的命题是这一陈述的更精确版本.
\setcounter{mtthm}{4}% 6.3、6.4 为步骤标题, 下一定理环境显示 6.5

\begin{proposition}[能量估计]\label{prop:6.5}
设 $\phi\in C([0,T],H^s)\cap C^1([0,T],H^{s-1})$ 为区间 $[0,T]$ 上 (67) 的解, 对某 $s\geq 1$ 满足 $\|\phi\|_{L_{[0,T]}^\infty L_x^\infty}\leq\Lambda_0$. 则 $\phi$ 满足如下能量估计:
\begin{equation}
\|\phi\|_{L_{[0,T]}^\infty\dot H^s}\leq C\bigl(\|\partial\phi\|_{L_{[0,T]}^1L_x^\infty},\Lambda_0\bigr)\|\phi[0]\|_{\dot H^s}.\end{equation}
此处常数 $C(a,b)$ 表示对 $a,b$ 的依赖, $\dot H^s$ 为通常的齐次 Sobolev 空间.\end{proposition}

\section*{6.6 步骤 2: 化归 Strichartz 型估计}

由 Cauchy--Schwarz 不等式, $\|\partial\phi\|_{L_{[0,T]}^1L_x^\infty}\leq T^{\frac12}\|\partial\phi\|_{L_{[0,T]}^2L_x^\infty}$. 于是 Strichartz 不等式 (71) 与时间区间 $[0,T]$ 之小性, 可用来在 $H^{s_*}$ 空间中闭合能量估计. 这就有效地给出了该问题所需的局部适定性. 定理 (A) 从而可化归如下自举论证.
\setcounter{mtthm}{6}% 6.6 为步骤标题, 下一定理环境显示 6.7

\begin{theorem}[A1]\label{thm:6.7}
设 $\phi\in C([0,T],H^{s_*})\cap C^1([0,T],H^{s_*-1})$ 为区间 $[0,T]$($T\leq 1$)上 (67) 的解. 假定
\begin{equation}
\|\phi\|_{L_{[0,T]}^\infty H^{s_*}}+\|\partial\phi\|_{L_{[0,T]}^2L_x^\infty}\leq B_0,\end{equation}
其中常数 $B_0\leq c_{s_*}^{-1}\Lambda_0$, $c_{s_*}$ 为嵌入 $H^{s_*}\subset L^\infty$ 的 Sobolev 常数. 则 $\phi$ 满足如下时间局部的 Strichartz 型估计:
\begin{equation}
\|\partial\phi\|_{L_{[0,T]}^2L_x^\infty}\leq C(B_0)T^{s_*-s_0}\|\phi\|_{L_{[0,T]}^\infty H^{s_*}}.\end{equation}
\end{theorem}

\setcounter{mtthm}{8}% 6.8 为步骤标题, 下一定理环境显示 6.9

\section*{6.8 步骤 3: Strichartz 型估计的二进版本与仿微分逼近}

为证明 Strichartz 型估计 (124), 我们可将方程 (67) 视为关于 $\phi$ 的、具粗糙系数的线性波动方程. 对系数作磨光, 并研究与一参数相关的、具光滑系数的线性波动方程族, 是有利的. 首先引入 $\phi^\lambda$: 将 $\phi$ 限制于傅里叶空间中频率 $\sim\lambda$ 的二进块得到. 更确切地, 设 $\zeta$ 为支在球壳 $\{\xi:\frac12\leq|\xi|\leq 2\}$ 上的光滑函数(此处 $\xi$ 表示空间傅里叶变换的变量), 并满足 $\sum_{k\in\mathbb{Z}}\zeta(2^k\xi)=1$($\forall\xi\in\mathbb{R}^n\setminus\{0\}$). 令 $\lambda$ 为二进参数 $\lambda=2^k$, $k\in\mathbb{Z}$, 记 $P_\lambda$ 为"投影子":
\begin{equation}
P_\lambda f(x)=f^\lambda(x)=\int e^{-\mathrm{i}x\cdot\xi}\zeta(\lambda^{-1}\xi)\hat f(\xi)\,d\xi.\end{equation}
又定义
\begin{equation}
f_{\leq\lambda}=S_\lambda f=\sum_{\mu\leq\lambda}P_\mu f.\end{equation}

定理 (A1) 由如下关于 $\phi^\lambda=P_\lambda\phi$ 的 Strichartz 型估计的二进版本推出.
\begin{theorem}[A2]\label{thm:6.9}
设 $\phi$ 如定理 (A1). 固定一大参数 $\Lambda$. 则对每个 $\lambda\geq\Lambda$, $\phi^\lambda$ 满足 Strichartz 型估计
\begin{equation}
\|\partial\phi^\lambda\|_{L_{[0,T]}^2L_x^\infty}\leq C(B_0)c_\lambda T^{s_*-s_0}\|\phi\|_{L_{[0,T]}^\infty H^{s_*}},\end{equation}
其中常数 $c_\lambda$ 满足 $\sum_\lambda c_\lambda^2\leq 1$. 类似估计对 $\phi_{\leq\Lambda}$ 亦成立.\end{theorem}

\textbf{注 (A2).} 对低频部分, 估计 (125) 对 $\phi_{\leq\Lambda}$ 由 Sobolev 不等式显然成立:
\begin{equation}
\|\partial\phi_{\leq\Lambda}\|_{L_{[0,T]}^2L_x^\infty}\leq cT^{\frac12}\|\phi_{\leq\Lambda}\|_{L_{[0,T]}^\infty H^{\frac52+\epsilon}}\leq c\Lambda^{\frac52+\epsilon-s_*}T^{\frac12}\|\phi\|_{L_{[0,T]}^\infty H^{s_*}},\end{equation}
其中 $c$ 为嵌入 $H^{\frac32+\epsilon}(\mathbb{R}^3)\subset L^\infty(\mathbb{R}^3)$ 的范数. 由于 $s_*$ 假定充分接近 $s_0=2+\frac12(2-\sqrt3)$, 且 $\Lambda$ 为可仅依赖于 $B_0$ 的固定大参数, 我们便得到 $\phi$ 低频部分所需的界.
我们只关注大频率 $\lambda\geq\Lambda$. 下一命题表明, 每个 $\phi^\lambda$ 满足一个具光滑度量 $g_{\leq\lambda^a}^{ij}=S_{\lambda^a}g^{ij}(S_\lambda^a\phi)$ 的非齐次波动方程, 其中参数 $a\in[0,1]$ 取固定值.
\begin{proposition}\label{prop:6.10}
设 $\phi$ 如定理 (A1). 固定参数 $a$ 的值, $a\in[0,1]$. 则对每个 $\lambda\geq\Lambda$, $\phi^\lambda$ 满足方程
\begin{equation}
\begin{aligned}
\Box_{g_{\leq\lambda^a}}\phi^\lambda&=-\partial_t^2\phi^\lambda+g_{\leq\lambda^a}^{ij}\partial_i\partial_j\phi^\lambda=R_\lambda^a,\\
\phi^\lambda|_{t=0}&=\phi_0^\lambda,\qquad \partial_t\phi^\lambda|_{t=0}=\phi_1^\lambda.\end{aligned}
\end{equation}
此外, 右端 $R_\lambda^a$ 的傅里叶支集含于集合 $\{\xi:\lambda\leq|\xi|\leq 4\lambda\}$ 中, 且对所有实数 $s>1$ 与任意 $t\in[0,T]$,
\begin{gather}
\|R_\lambda(t)\|_{\dot H^s}\leq c\lambda^s\|R_\lambda\|_{L_x^2}\lesssim C(B_0)\lambda^{1-a}c_\lambda\|\partial\phi\|_{L_x^\infty}\|\phi\|_{\dot H^s},\\
c_\lambda:\ \sum_\lambda c_\lambda^2\leq 1.\end{gather}
\end{proposition}

\setcounter{mtthm}{11}% 6.11 为步骤标题, 下一定理环境显示 6.12

\section*{6.11 步骤 4: 频率相关区间上的 Strichartz 估计}

此步骤将定理 (A) 的证明化归为如下线性化方程 $\Box_{g_{\lambda^a}}\psi=0$ 在小时间区间 $I$(大小 $\approx T\lambda^{-(1-a)}$)上精确 Strichartz 估计的证明. 最终 Strichartz 估计中正则性的损失, 正是将 $[0,T]$ 上这些尖锐的 Strichartz 估计对区间 $I$ 求和的结果.
\begin{theorem}[A3]\label{thm:6.12}
设 $\lambda\geq\Lambda$, $\psi$ 为线性波动方程 $\Box_{g_{\lambda^a}}\psi=0$ 的解, 初数据 $\psi[0]$ 满足 $\operatorname{supp}\hat\psi[0]\in\{\frac12\lambda\leq|\xi|\leq 2\lambda\}$, 度量 $g_{\lambda^a}$ 如\propref{prop:6.10} 所定义. 固定参数 $a$ 的值, $a=\frac{5-2s_0}{1-2s_0+2s_*}$ 且 $a<-1+\sqrt3$. 则存在区间 $[0,T]$ 的一个分拆 $\{I\}$ 为子区间 $I$, 使得每个 $I$ 的大小 $|I|\leq T\lambda^{-(1-a)}$, 子区间的个数约为 $\lambda^{1-a}$, 且在每个 $I$ 上, $\psi$ 满足一带任意固定小量 $\epsilon>0$ 的 Strichartz 估计:
\begin{equation}
\|P_\lambda\partial\psi\|_{L_I^2L_x^\infty}\leq C(B_0)|I|^\epsilon\|\psi[0]\|_{\dot H^{2+\epsilon}}.\end{equation}
\end{theorem}

\textbf{注 (A3).} 依自举假设 (123), 拟线性问题的解 $\phi$ 满足估计 $\|\partial\phi\|_{L_{[0,T]}^2L_x^\infty}\leq B_0$. 由此易见存在时间区间 $[0,T]$ 的一个子分拆 $\{I\}$, 其总子区间数介于 $\lambda^{1-a}$ 与 $2\lambda^{1-a}$ 之间, 每个 $I$ 的大小不超过 $T\lambda^{-(1-a)}$, 且在每个 $I$ 上我们有
\begin{equation}
\|\partial\phi\|_{L_I^2L_x^\infty}\leq\lambda^{-\frac{1-a}{2}}\|\partial\phi\|_{L_{[0,T]}^2L_x^\infty}.\end{equation}
这一构造定义了定理 (A3) 中提到的子分拆 $\{I\}$.
\section*{6.13 度量 $g_{\lambda^a}$ 的性质}

不等式 (128) 是具变系数(度量)波动方程 $\Box_{g_{\leq\lambda^a}}\psi=0$ 在频率相关区间 $I$ 上的 Strichartz 估计.
度量 $g_{\leq\lambda^a}=S_{\lambda^a}g(S_{\lambda^a}\phi)$ 依赖于拟线性问题的解 $\phi$. 下一命题陈述了由 $\phi$ 上的自举条件 (123) 与注 (A3) 中区间分拆的构造所导出的族 $g_{\leq\lambda^a}$ 的性质.
\setcounter{mtthm}{13}% 6.13 为步骤标题, 下一定理环境显示 6.14

\begin{proposition}\label{prop:6.14}
设 $\phi\in C([0,T],H^{s_*})\cap C^1([0,T],H^{s_*-1})$ 为区间 $[0,T]$($T\leq 1$)上 (67) 的解. 假设 $\phi$ 满足定理 (A1) 的假设 (123). 考虑由注 (A3) 定义的子分拆 $\{I\}$. 则度量族 $g_{\leq\lambda^a}=S_{\lambda^a}g(S_{\lambda^a}\phi)$ 满足如下条件:
对所有子区间 $I$ 与所有非负整数 $m$,
\begin{gather}
\|\partial^{1+m}g_{\leq\lambda^a}\|_{L_I^1L_x^\infty}\leq\lambda^{-(1-a)+am}\bar B_0,\\
\|\partial^{1+m}g_{\leq\lambda^a}\|_{L_I^2L_x^\infty}\leq\lambda^{-\frac{1-a}{2}+am}\bar B_0,\\
\|\partial^{1+m}g_{\leq\lambda^a}\|_{L_I^\infty L_x^\infty}\leq\lambda^{\frac{a^2}{2}+am}\bar B_0,\\
\|\partial_x^{\frac12+m}(\partial^2 g_{\leq\lambda^a})\|_{L_I^\infty L_x^2}\leq\lambda^{\frac{a^2}{2}+am}\bar B_0,\\
\|\partial^m\Box_{g_{\leq\lambda^a}}g_{\leq\lambda^a}\|_{L_I^1L_{\bar x}^\infty}\leq\lambda^{-(1-a)+am}\bar B_0.\end{gather}
常数 $\bar B_0$ 只依赖于常数 $M_0$ 与 $B_0$.\end{proposition}

注: 由构造可知, 度量 $g_{\leq\lambda^a}$ 的频率只在相对空间变量 $x$ 的傅里叶对偶变量方向被截断于 $\lambda^a$ 之上. 因此, 关于 $x$ 的每一次微分至多引入一个因子 $\lambda^a$. 然而, 利用 $g_{\leq\lambda^a}$ 依赖于波动方程之解这一事实, 对时间导数我们也能得到同样的结论.
定理 (A3) 可重新表述为: 关于具背景度量 $g_{\leq\lambda^a}$(满足\propref{prop:6.14} 之估计)的线性波动方程, 在固定小区间 $I$ 上局部 Strichartz 估计的结果.
\begin{theorem}[A4]\label{thm:6.15}
设 $\psi$ 为线性波动方程
\begin{equation}
\begin{aligned}
\Box_{g_{\leq\lambda^a}}\psi&=-\partial_t^2\psi+g_{\leq\lambda^a}^{ij}\partial_i\partial_j\psi=0,\\
\psi|_{t=0}&=\psi_0,\qquad \partial_t\psi|_{t=0}=\psi_1
\end{aligned}
\end{equation}
在区间 $I$(长度 $|I|\leq\lambda^{-(1-a)}$)上的解, 初数据 $\psi[0]$ 在傅里叶空间中支于集合 $\{\xi:\frac12\lambda\leq|\xi|\leq 2\lambda\}$. 假设度量 $g_{\lambda^a}$ 满足\propref{prop:6.14} 的 (130)--(134), 参数 $a$ 取 $a<-1+\sqrt3$. 令 $P_\lambda$ 为到傅里叶空间集合 $\{\xi:\frac12\lambda\leq|\xi|\leq 2\lambda\}$ 的投影. 则对充分大参数 $\Lambda$、所有二进 $\lambda\geq\Lambda$ 与固定 $\epsilon>0$,
\begin{equation}
\|P_\lambda\partial\psi\|_{L_I^2L_x^\infty}\leq C(\bar B_0)|I|^\epsilon\|\psi[0]\|_{\dot H^{2+\epsilon}},\end{equation}
其中常数 $C(\bar B_0)$ 与 $\lambda$ 无关.\end{theorem}

\setcounter{mtthm}{16}% 6.16 为步骤标题, 下一定理环境显示 6.17

\section*{6.16 步骤 5: 重标度}

将问题 (135) 替换为它的重标度版本是方便的, 使得初数据的频率支集满足 $|\xi|\sim 1$, 且时间区间 $I$ 的长度为 $\leq\lambda^a$.
必要时对问题作时间平移, 可设时间区间 $I$ 始于 $t=0$. 引入重标度度量族 $h_\lambda$:
\begin{equation}
h_\lambda(t,x)=g_{\leq\lambda^a}(\lambda^{-1}t,\lambda^{-1}x).\end{equation}
\propref{prop:6.14} 推出, $h_\lambda$ 在时间区间 $I=[0,t_*]$($t_*\leq\lambda^a$)上满足如下估计:
\begin{gather}
\|\partial^{1+m}h_\lambda\|_{L_{I_\lambda}^1L_x^\infty}\lesssim\lambda^{-(1-a)(m+1)},\\
\|\partial^{1+m}h_\lambda\|_{L_{I_\lambda}^2L_x^\infty}\lesssim\lambda^{-\frac{2-a}{2}-(1-a)m},\\
\|\partial^{1+m}h_\lambda\|_{L_{I_\lambda}^\infty L_x^\infty}\lesssim\lambda^{-1+\frac{a^2}{2}-(1-a)m},\\
\|\partial_x^{\frac12+m}(\partial^2 h_\lambda)\|_{L_{I_\lambda}^\infty L_x^2}\lesssim\lambda^{-1+\frac{a^2}{2}-(1-a)m},\\
\|\partial^m\Box_{h_\lambda}h_\lambda\|_{L_{I_\lambda}^1L_x^\infty}\lesssim\lambda^{-(2-a)-(1-a)m}.\end{gather}

重标度后, 定理 (A4) 化为

\begin{theorem}[A5]\label{thm:6.17}
设 $\psi$ 为线性波动方程
\begin{equation}
\begin{aligned}
\Box_{h_\lambda}\psi&=-\partial_t^2\psi+h_\lambda^{ij}\partial_i\partial_j\psi=0,\\
\psi|_{t=0}&=\psi_0,\qquad \partial_t\psi|_{t=0}=\psi_1
\end{aligned}
\end{equation}
在区间 $[0,t_*]$($t_*\leq\lambda^a$)上的解. 假设参数 $\lambda\geq\Lambda$($\Lambda$ 为充分大常数), 且度量 $h_\lambda$ 满足 (138)--(142), 参数 $a$ 满足 $a<-1+\sqrt3$. 令 $P$ 为到傅里叶空间集合 $\{\xi:1\leq|\xi|\leq 2\}$ 的投影算子. 则
\begin{equation}
\|P\partial\psi\|_{L_{[0,t_*]}^2L_x^\infty}\lesssim|t_*|^\epsilon\bigl(\|\partial\psi_0\|_{L_x^2}+\|\psi_1\|_{L_x^2}\bigr),\end{equation}
其中不等式 $\lesssim$ 中的常数与 $\lambda$ 无关.\end{theorem}

注: 定理 (A5) 不含关于初数据 $\psi[0]$ 傅里叶支集的任何假设.
\section*{6.18 步骤 6: 衰减估计}

标准 $TT^*$ 型论证的一个变体(见 \cite{klai29})使我们能将 Strichartz 估计 (144) 化归为相应的色散不等式, 即 (145). 在此过程中, 我们将方程 $\Box_{h_\lambda}\psi=0$ 替换为几何波动方程
\begin{equation}
\Box_{h_\lambda}\psi=-\frac{1}{\sqrt{\det h_\lambda}}\partial_t\sqrt{\det h_\lambda}\,\partial_t\psi+\frac{1}{\sqrt{\det h_\lambda}}\partial_i\bigl(h_\lambda^{ij}\sqrt{\det h_\lambda}\,\partial_j\psi\bigr)=0.\end{equation}

\setcounter{mtthm}{18}% 6.18 为步骤标题, 下一定理环境显示 6.19

\begin{theorem}[A6]\label{thm:6.19}
设 $\psi$ 为线性波动方程
\begin{equation}
\begin{aligned}
\Box_{h_\lambda}\psi&=-\frac{1}{\sqrt{\det h_\lambda}}\partial_t\sqrt{\det h_\lambda}\,\partial_t\psi+\frac{1}{\sqrt{\det h_\lambda}}\partial_i\bigl(h^{ij}_\lambda\sqrt{\det h_\lambda}\,\partial_j\psi\bigr)=0,\\
\psi|_{t=0}&=\psi_0,\qquad \partial_t\psi|_{t=0}=\psi_1
\end{aligned}
\end{equation}
在区间 $[0,t_*]$($t_*\leq\lambda^a$)上的解, 初数据 $\psi[0]$ 在傅里叶空间中支于集合 $\{\xi:\frac12\leq|\xi|\leq 2\}$. 我们只考虑大参数值 $\lambda\geq\Lambda$. 假设度量 $h_\lambda$ 满足 (138)--(142), 参数 $a$ 满足 $a<-1+\sqrt3$. 则对所有 $t\leq t_*$ 与任意固定小量 $\epsilon>0$,
\begin{equation}
\|P\partial\psi(t)\|_{L_x^\infty}\lesssim\frac{1}{(1+|t|)^{1-\epsilon}}\|\psi[0]\|_{L_x^1}.\end{equation}
\end{theorem}

我们通过在物理空间将初数据 $\psi[0]$ 分解为支集本质不交、含于半径 $\frac12$ 的球中的函数之和, 作最后的化简. 利用 $L^1$ 范数的可加性与标准 Sobolev 不等式, 我们可将色散不等式 (146) 化归为 $L^2$--$L^\infty$ 衰减估计.
\begin{theorem}[B]\label{thm:6.20}
设 $\psi$ 为线性波动方程 (145) 在区间 $[0,t_*]$($t_*\leq\lambda^a$)上的解, 初数据 $\psi[0]$ 支于物理空间中以原点为中心、半径 $\frac12$ 的球 $B_{\frac12}(0)$ 内. 我们固定一大常数 $\Lambda$, 并只考虑大参数值 $\lambda\geq\Lambda$. 假设度量 $h_\lambda$ 满足 (138)--(142), 参数 $a$ 满足 $a<-1+\sqrt3$. 则对所有 $t\leq t_*$、任意小量 $\epsilon>0$ 与充分大整数 $m>0$,
\begin{equation}
\|P\partial\psi(t)\|_{L_x^\infty}\lesssim\frac{1}{(1+|t|)^{1-\epsilon}}\sum_{k=1}^m\|\partial^k\psi[0]\|_{L_x^2}.\end{equation}
\end{theorem}
